\chapter{Qualité, validation et évolutions}
\label{chap:qualite}

\section{État de validation}

\subsection{Périmètre de lecture au 22 octobre 2026}

Le master courant est Rec.~2020 linéaire D65. La calibration RAW est active par défaut, sous réserve de compatibilité des règles du catalogue avec le boîtier, le profil et le mode. Le sidecar courant est v6 et mémorise le développement RAW et les réglages optiques. Les caches partagés et le contrôleur mémoire sont intégrés ; leur qualification sur tous les systèmes et corpus reste distincte de leur présence dans le code.

La matrice de validation et les JSON de campagne font référence pour le statut de chaque groupe. Les chiffres datés ci-dessous décrivent des campagnes historiques, pas une nouvelle exécution au 22 octobre. Les travaux non fusionnés de \texttt{patch\_color}, \texttt{patch\_noise} et \texttt{patch\_ci\_cd} ne sont pas présumés livrés dans \texttt{dev}.

\subsection{Campagnes historiques d'octobre 2026}

Les ajustements fondamentaux ont été comparés à des exports de référence sur un corpus TIFF en sRGB : exposition, noirs, ombres, blancs, hautes lumières, contraste et température. Les mesures utilisent notamment les différences par pixel (MAE) et des inspections de superposition lorsque les dimensions et l'encodage le permettent. Cette méthode distingue un décalage spatial d'une différence de rendu.

Le mélangeur huit bandes dispose de campagnes aux extrêmes et aux valeurs
intermédiaires. Le bruit v2 a une progressivité mesurée sur douze scènes ;
cela ne garantit pas la conservation de tous les détails. Color grading,
courbes, netteté, grain/flou et leurs interactions nécessitent une recette
photographique. Sidecars v4 et application des réglages à une sélection étaient alors
implémentés ; leur recette manuelle complète reste ouverte. L'export par lot
n'est pas implémenté. Le DNG de rendu n'est pas certifié compatible Lightroom.

Au 8 octobre, Camera Neutral est accepté visuellement par l'utilisateur.
La normalisation était alors expérimentale, limitée à cinq fichiers exacts et
désactivée par défaut ; cette restriction n'est plus le contrat courant. La suite complète WSL passe 17 tests sur 18 : la sonde
JPEG 9933 reste en échec (MAE RGB 25,1519 pour un seuil 25), hors mode calibré.
Les cinq suites Windows ciblées de la dernière adaptation portable et les
76 tests Python d'analyse plus 12 de sélection réussissent. Ces nombres sont
une photographie datée, pas une garantie d'exécution future.

\section{Éléments à compléter avant une distribution}

\begin{itemize}
  \item Établir des jeux de tests reproductibles pour chaque réglage et enregistrer les métriques dans un format structuré.
  \item Mesurer la durée d'ouverture sur des RAW représentatifs, le pic mémoire et l'efficacité réelle du cache LRU et du préchargement de vignettes.
  \item Formaliser la compatibilité des profils DCP/LCP : origine, licence, version et droit de redistribution de chaque fichier.
  \item Compléter la recette du schéma sidecar v6 et de la persistance RAW/optique déjà implémentée ; vérifier les conflits, les déplacements et les ressources manquantes.
  \item Ajouter un test de non-régression pour le rendu des profils, le rendu en lot et les métadonnées.
  \item Documenter le format d'export final (profondeur, profil ICC incorporé, métadonnées EXIF/XMP) et ses conditions d'utilisation.
\end{itemize}

\section{Gouvernance documentaire appliquée}

Les rapports expérimentaux sont répartis entre \texttt{docs/validation} (protocole, analyse et mesures) et \texttt{docs/reports} (journal de séance). La présente arborescence est la documentation normative compilable. Toute fonctionnalité terminée doit mettre à jour son chapitre, ses diagrammes concernés, les références bibliographiques, la matrice de validation et un scénario de test. Les décisions d'architecture transverses doivent être ajoutées à une courte ADR (\emph{Architecture Decision Record}) afin de ne pas dépendre exclusivement des rapports de session.

\section{Mesures exportées et graphiques reproductibles}

\texttt{scripts/export\_documentation\_metrics.py} génère les CSV par famille
et les graphiques SVG depuis les preuves immuables. Chaque valeur conserve
son fichier source, son empreinte SHA-256 et son emplacement JSON ou textuel.
Le répertoire \texttt{docs/validation/metrics} est une vue dérivée, pas une
nouvelle source de vérité. Les nombres extraits littéralement d'un journal
ne deviennent pas des métriques homogènes sans interprétation de leur contexte.

Les durées CPU et les indicateurs de rendu sont séparés : la courbe de temps
représente seulement le benchmark Windows historique avant/après parallélisation.
Les réponses tonales par révision, la progression du débruitage et les écarts
DCP à Lightroom ne sont pas des mesures de vitesse. Aucun temps actuel de
rendu, de vignettes ou d'export n'est inventé ; des campagnes restent nécessaires.
